\section{构建记忆层：实现要点}

原文以 Python + OpenAI Embeddings + ChromaDB 为例，展示了如何构建一个记忆层。核心组件包括两个类。

\subsection{MemoryStore 类}

负责记忆的写入（生成 Embedding 后存储）和语义检索，是整个系统的基础。

\begin{lstlisting}[caption={MemoryStore 核心接口（伪代码）}]
class MemoryStore:
    def __init__(self, collection_name):
        self.client = chromadb.Client()
        self.collection = self.client.get_or_create_collection(
            collection_name)
        self.embedder = OpenAI()

    def add(self, text, metadata=None):
        embedding = self.embedder.embed(text)
        self.collection.add(embedding, text, metadata)

    def search(self, query, top_k=5):
        query_embedding = self.embedder.embed(query)
        return self.collection.query(query_embedding, n=top_k)
\end{lstlisting}

\subsection{EpisodicLogger 类}

在 MemoryStore 之上添加情景日志层，记录每次任务的执行过程和结果。

\begin{lstlisting}[caption={EpisodicLogger 核心接口（伪代码）}]
class EpisodicLogger:
    def __init__(self, memory_store):
        self.store = memory_store

    def log_episode(self, task, strategy, outcome, lesson):
        episode = {
            "task": task, "strategy": strategy,
            "outcome": outcome, "lesson": lesson,
            "timestamp": now()
        }
        self.store.add(json.dumps(episode),
                       metadata={"type": "episode"})

    def recall_similar(self, task_description, top_k=3):
        return self.store.search(task_description, top_k)
\end{lstlisting}

\subsection{本章小结}

记忆层的核心是 MemoryStore（向量存储 + 语义检索）和 EpisodicLogger（行动日志 + 经验回溯）。两者搭配即可构建一个具备基本记忆能力的 Agent。

